\chapter{Metodologia}
O AADL é empregado desde o início do processo do desenvolvimento, decisões arquiteturais são feitas após extensa análise para garantir a melhor eficiência e confiabilidade. Os modelos dos componentes são gerados e refinados iterativamente.
No início do processo de modelagem, o fluxo da comunicação, os dispositivos e os comportamentos são definidos. Então, os componentes são descritos com AADL. Isto inclui sensores, atuadores, unidades de processamento e o software que orquestra as funções finais do sistema. 

Assim que os componentes são identificados, uma hierarquia modular de pacotes é criada, sempre visando reutilização e clareza. Estes pacotes representam agrupamentos de componentes, como lógica de controles, interfaces de sensores e subsistemas de mobilidade. Na modelagem de Robinho, o sistema de direção diferencial e o ESP-Cam são modelados como dispositivos separados, enquanto a unidade de processamento principal é definida como um sistema com processos que integram informações de sensores e geram comandos para o motor.

Após a definição da estrutura lógica, os próximos passos envolvem alimentar o modelo com propriedades relacionadas a performance, como tempo de execução de threads, latência de comunicação, frequências, e propriedades de escalonamento de processos. Estas propriedades permitem uma análise formal de agendamento e contagem de tempo usando o OSATE.

Neste ponto, o modelo é instanciado. Isto permite checagem de validação para consistência arquitetural, corretude de fluxo de dados e restrições temporais. No caso de Robinho, isto inclui verificar se o dado consegue ser processado e transmitido dentro do tempo limite para gerar comandos de movimento com latência mínima e máxima eficiência.

Finalmente, depois de o modelo ser instanciado, o OSATE permite o uso das ferramentas de análise integradas, como análise de latência, simulação de escalonamento e previsão de uso de recursos, para definir se a arquitetura atual cumpre os requisitos do sistema. Caso quaisquer gargalos ou violações sejam detectadas, o modelo pode ser refinado iterativamente, suportando uma metodologia de design antes da implementação, que reduz significativamente os riscos de integração durante a implantação.

Repetidas vezes durante o processo de modelagem, foram encontrados gargalos computacionais e na comunicação entre os componentes. Nestes casos, O uso dos componentes específicos e meios de comunicação foram reavaliado, e iterativamente os componentes do sistema foram reavaliados, e se necessário, alterados.

No caso de a aplicação final do sistema alterar, por exemplo, os módulos de software forem trocados, a única medida necessária é que os componentes novos sejam modelados e testados. Isto torna o sistema dos modelos altamente escalável e interoperável, sem necessidade de retrabalho modelando novamente os dispositivos.

Esta abordagem dirigida a modelagem garante que as decisões arquiteturais por trás dos sistemas sejam documentadas e validadas contra métricas quantitativas de performance muito cedo no ciclo de desenvolvimento.

Ao finalizar todos os processos de análise, é possível implantar o sistema em dispositivos reais com significativamente menos probabilidade de riscos como gargalos ou falta de poder computacional.

